\documentclass{thecours}
\begin{document}
\bigpart{Additions et soustractions de nombres relatifs}
\methode{additions_soustractions_decimaux}
\par\noindent{\bfseries Exemple :} La somme des termes 12,67 et 29,8.\vskip5pt
\par\noindent{\bfseries Exemple :} La différence des termes 29,8 et 12,67.
\bigpart{Multiplication}
\subpart{Calcul posé}
\methode{multiplication_decimaux}
\par\noindent{\bfseries Exemple :} Le produit de 29,84 et 1,7.\vskip15pt
\subpart{Calcul mental}
\methode{multiplication_par_05_025}
\par\noindent{\bfseries Exemple : }\vskip15pt
\hbox{\kern30pt\hbox to.5\linewidth{$68\times0,5=34$ car $34\times2=68$\hfill}%
\hbox{$0,5\times15=7,5$ car $7,5\times2$\hfill}}\vskip10pt
\hbox{\kern30pt\hbox to.5\linewidth{$84\times0,25=21$ car $21\time4=84$\hfill}%
\hbox{$0,25\times18=4,5$ car $4,5\times4=18$\hfill}}\vskip10pt
\bigskip\methode{multiplier_par_01_001}
\par\noindent{\bfseries Exemple : }\vskip15pt
\hbox{\kern30pt\hbox to.5\linewidth{$120\times0,1=12$\hfill}%
\hbox{$0,1\times7,2=0,72$\hfill}}\vskip10pt
\hbox{\kern30pt\hbox to.5\linewidth{$847,6\times0,01=8,476$\hfill}%
\hbox{$0,001\times34=0,034$\hfill}}\vskip10pt
\bigpart{Division décimale}
\methode{division_decimale}\vskip10pt
\noindent{\bfseries Exemple :}\vskip10pt\par\noindent\kern50pt
\opdiv[displayintermediary=nonzero,shiftdecimalsep=none]{12.4}{5}\kern50pt
Donc $12,4=5\times2,48$.
\end{document}
